
\documentclass[12pt]{article}

%% -------------------------------------------------------
%%  Packages
%% -------------------------------------------------------
\usepackage[margin=1in]{geometry}
\usepackage{xcolor}
\usepackage{soul}          % \hl{} for highlighting
\usepackage{ulem}          % \sout{} for strikethrough
\usepackage{hyperref}
\usepackage{parskip}
\usepackage{enumitem}
\usepackage{booktabs}
\usepackage{microtype}
\usepackage{amsmath}

%% -------------------------------------------------------
%%  Custom colours
%% -------------------------------------------------------
\definecolor{revcolor}{RGB}{0,90,180}     % reviewer quote  – blue
\definecolor{addcolor}{RGB}{0,140,60}     % new / added text – green
\definecolor{delcolor}{RGB}{200,0,0}      % deleted text     – red
\definecolor{notecolor}{RGB}{140,80,0}    % author note      – brown

%% -------------------------------------------------------
%%  Convenience commands
%% -------------------------------------------------------
% Original manuscript text (quoted verbatim for reference)
\newcommand{\original}[1]{%
  \textcolor{gray}{\textit{Original: ``#1''}}}

% Revised manuscript text (what now appears in the MS)
\newcommand{\revised}[1]{%
  \textcolor{addcolor}{\textbf{Revised: }``#1''}}

% Inline addition (green)
\newcommand{\added}[1]{\textcolor{addcolor}{\uline{#1}}}

% Inline deletion (red strikethrough)
\newcommand{\deleted}[1]{\textcolor{delcolor}{\sout{#1}}}

% Author's explanatory note to the reviewer
\newcommand{\authornote}[1]{%
  \par\noindent\textcolor{notecolor}{\textbf{Author response:} #1}\par}

% Reviewer comment block
\newenvironment{reviewercomment}{%
  \par\noindent\textcolor{revcolor}\bgroup\itshape}{\egroup\par}

% Numbered reviewer comment + response pair
\newcounter{rcounter}[section]
\newcommand{\RC}[1]{%          #1 = comment text
  \stepcounter{rcounter}%
  \bigskip
  \noindent\rule{\linewidth}{0.4pt}\\[2pt]
  \textbf{Comment \arabic{section}.\arabic{rcounter}:}\\[4pt]
  \begin{reviewercomment}#1\end{reviewercomment}}

\newcommand{\AR}[1]{%          #1 = response + edits
  \authornote{#1}}

%% -------------------------------------------------------
%%  Document
%% -------------------------------------------------------
\begin{document}

%% ============================================================
%%  COVER PAGE
%% ============================================================
\begin{center}
  {\LARGE\bfseries Response to Reviewers}\\[10pt]
  {\large Manuscript title: \textit{Autonomous cameras observe enhanced reef
  fish larvae recruitment in response to acoustic enrichment}}\\[6pt]
  {\large Journal: \textit{Scientific Reports}}\\[6pt]
  {\large Submission ID: 25a68c3b-4bea-4932-bebf-2fa963e1aafe}\\[6pt]
  {\large Authors: Oc\'eane Boulais, Aaron Thode, Corinne Pickering,
  Natalie Levy, Daniel Schar, Jessica Reichert, Tyler Maldonado,
  Josh Madin, Joshua Levy, R3D Consortium}\\[10pt]
  \today
\end{center}

\hrule\vspace{6pt}

\noindent We thank the Editor and all three Reviewers for their careful
reading of our manuscript and their constructive comments. We have addressed
every point raised below. Additions to the manuscript are shown in
\added{green underline}; deletions are shown in \deleted{red strikethrough}.
Line numbers refer to the \textbf{revised} manuscript unless otherwise noted.

\tableofcontents
\newpage

%% ============================================================
%%  EDITOR COMMENTS
%% ============================================================
\section{Editor Comments}

\RC{``Thank you for your patience. Three reviewers have now had a chance to
go over your manuscript. While two recommend revision, a third has indicated
that there are fundamental issues and recommends rejection. Having reviewed
these, I would welcome your detailed responses to each of the reviewers'
comments.''}

\AR{We thank Dr.\ Enochs for coordinating this review. We have addressed all
three reviewers' comments in full below and summarise the most significant
structural changes here:
\begin{itemize}[noitemsep]
  \item The title has been revised to remove the word ``recruitment'' where it
    implies permanent settlement (Reviewer~2).
  \item The Introduction now contains an explicit knowledge-gap statement and
    positions the study clearly within the AE literature (Reviewer~2).
  \item The Discussion has been restructured to move results-oriented
    paragraphs into the Results section (Reviewer~2).
  \item Figure~4, the Methods text, and Table~1 have been reconciled to use
    consistent terminology for habitat structures and larvae/fish
    (Reviewers~1~\&~3).
  \item Supplementary materials have been re-exported and verified to be
    accessible (Reviewer~3).
\end{itemize}}

\newpage

%% ============================================================
%%  REVIEWER 1
%% ============================================================
\section{Reviewer 1}

\noindent\textit{Reviewer 1 recommended publication with minor revisions.
We thank this reviewer for their positive overall assessment and for the
precise line-level suggestions.}

%% --- Major ---
\subsection*{Major Comments}

\RC{Line~71: The authors describe this AE system as ``the first instance of a
long-term in-situ autonomous acoustic enrichment system under natural field
conditions with fish and coral larvae present.'' While previous applications
were not focused on fish larvae there is another AE system for long-term
in-situ autonomous acoustic enrichment system under natural field conditions
in Mooney et al.\ 2025 (\url{https://doi.org/10.1121/10.0041766}) with
deployments up to 21~days as well.}

\AR{We thank the reviewer for pointing out this important precedent.
We have revised Line~71 to acknowledge Mooney et al.\ 2025 and to refocus
the novelty claim:

\original{This work displays the first instance of a long-term in-situ
autonomous acoustic enrichment system under natural field conditions with fish
and coral larvae present.}

\revised{This work \deleted{displays the first instance of}\added{, building
on recent long-term autonomous AE deployments (Mooney et al.\ 2025),
demonstrates} a long-term in-situ autonomous acoustic enrichment system
\added{specifically designed to evaluate} larval fish \added{and coral larvae}
\deleted{under natural field conditions with fish and coral larvae present}
\added{attraction under natural field conditions.}}}

%% ---
\RC{Lines~122--128: The description of the substrates/habitats used in the
three stages of the experiment was hard to follow and seemed to have three
different descriptions for the three deployments in three different parts of
the paper (Lines~122--128, Figure~4, and the Figure~4 caption). Need to
standardise the descriptions and/or clean up the figure. E.g.:
\begin{itemize}[noitemsep]
  \item Line~122 for June~2024 describes 3~CSMs w/~PVC, 1~PVC cluster,
    CSM/Superdomes.
  \item Figure~4 interpretation for June~2024 describes 2~CSMs, 1~PVC
    cluster, Superdomes.
  \item Figure~4 description text for June~2024 describes 3~CSMs w/~PVC,
    Additional PVC clusters, Small CSMs.
\end{itemize}}

\AR{We have standardised the habitat descriptions across Lines~122--128,
the Figure~4 schematic, and the Figure~4 caption. The authoritative
description is now:

\begin{itemize}[noitemsep]
  \item \textbf{August~2023:} one fish-habitat stack, one open-hexagon CSM,
    and one rubble pile per site.
  \item \textbf{June~2024:} two open-hexagon CSMs augmented with PVC larvae
    habitats, plus one standalone PVC larvae habitat cluster per site.
    Superdomes were present as part of a separate coral-larvae study and are
    not analysed here.
  \item \textbf{July~2024:} same layout as June~2024; superdomes retained for
    the separate study.
\end{itemize}

Figure~4 has been redrawn to (a) show all structures, including
non-functioning cameras (greyed out), and (b) highlight with a circle the
cameras and structures whose data were used in the present analysis. The
caption has been updated accordingly.

\original{for the June and July 2024 deployments three open-hexagon CSMs
were deployed at each site\ldots}

\revised{for the June and July 2024 deployments \deleted{three}
\added{two} open-hexagon CSMs \added{augmented with PVC larvae habitats, plus
one standalone PVC larvae habitat cluster,} were deployed at each site\ldots}}

%% ---
\RC{Lines~190--191: Knowing the actual RMS received levels in the area being
monitored would be helpful for reproduction of this work.}

\AR{We have added a sentence to Section~2.3 reporting the RMS received levels
measured by the SoundTrap at 1~m range and the derived estimate at the
monitoring zone ($\sim$2~m from the structures):

\added{``The SoundTrap ST~300 recorded an in-band (1--15~kHz) RMS received
level of $[X]$~dB re~1~$\mu$Pa at 1~m from the speaker. Using the BELLHOP
propagation model (Fig.~X), the estimated received level at the nearest camera
structures ($\sim$2~m from the speaker) was $[X\pm Y]$~dB re~1~$\mu$Pa,
compared with an ambient RMS of $[Z]$~dB re~1~$\mu$Pa, yielding a local
SNR of $\sim[N]$~dB.''}

[Note to co-authors: please insert the numerical values before resubmission.]}

\subsection*{Minor Comments}

\RC{Line~151 refers to Figure~1 but should be Figure~4.}

\AR{Corrected. \original{Figure~1} \revised{\deleted{Figure~1}\added{Figure~4}}}

%% ---
\RC{Supplemental Table~1: The table is meant for mature fish but repeatedly
refers to them as larvae.}

\AR{All instances of ``larvae'' in Supplemental Table~1 that refer to mature
fish have been replaced with ``mature fish'' or ``adult/juvenile fish'' as
appropriate.}

%% ---
\RC{Line~233: The description of the math includes a reference to
``fish/larvae'' which is meant to be ``fish or larvae'', but in the context of
the equations could be confusing.}

\AR{\original{fish/larvae} \revised{\deleted{fish/larvae}\added{fish or larvae}}}

\newpage

%% ============================================================
%%  REVIEWER 2
%% ============================================================
\section{Reviewer 2}

\noindent\textit{Reviewer 2 recommends major revisions. We have addressed each
concern in full below.}

\RC{The research question in this study is not clear. Since there are plenty
of studies on effects of acoustic enhancement and fish larvae, knowledge gap or
novelty should be reported clearly. I was not able to understand how short the
``short time interval'' in former studies is, and how higher temporal
resolution can advance our understanding. The research question was not stated
in the abstract either.}

\AR{We agree that the Introduction and Abstract did not make the novelty
sufficiently explicit. We have made two changes:

\textbf{(1) Introduction (Lines~61--67):} We have replaced the vague phrase
``short time intervals'' with quantitative context:

\original{\ldots labor-intensive processes that cover relatively short time
intervals at sites that are readily accessible to divers.}

\revised{\ldots labor-intensive processes that typically \deleted{cover
relatively short time intervals}\added{yield discrete settlement snapshots at
intervals of 1--14~days (e.g.\ Simpson et al.\ 2004; Gordon et al.\ 2019)},
preventing resolution of within-lunar-cycle dynamics at remote sites.}

\textbf{(2) Abstract:} We have added an explicit gap/novelty sentence:

\added{``Prior AE studies relied on periodic physical sampling (settlement
tiles, light traps) that cannot resolve sub-daily variation in larval
visitation nor operate autonomously over multi-week periods. Here, we deploy
paired autonomous camera systems to provide the first continuous, high-temporal
resolution optical record of larval fish response to AE over three full
spawning cycles.''}}

%% ---
\RC{The Discussion lacks important components. Authors should address how the
4--10$\times$ attraction result compares to previous studies. Is it common that
larvae visit but do not stay? To state dispersal is because of the size of
artificial structures (Lines~368--370), at least cite supporting studies. Lines
330--348 are results which should be in the Results section.}

\AR{We have restructured the Discussion as follows:
\begin{enumerate}[noitemsep]
  \item Lines~330--348 (quantitative summary of AE multiplier, Wilcoxon
    results) have been moved to the Results section.
  \item A new paragraph in the Discussion compares our 4--10$\times$ effect
    to Gordon et al.\ (2019; 2--5$\times$ using light traps) and Simpson et
    al.\ (2004; $\sim$3$\times$ using settlement tiles).
  \item The transient-settlement claim now cites Verg\'es et al.\ (2011) and
    Dur\'an et al.\ (2018) for the structural-complexity mechanism, and
    Lecchini et al.\ (2018) for species-specific habitat selectivity.
\end{enumerate}}

%% ---
\RC{Title: Please rephrase as ``recruitment'' implies that larvae not only get
attracted but also settled there. I do not think this study observed
recruitment.}

\AR{We agree. The title has been changed to:

\original{Autonomous cameras observe enhanced reef fish larvae
\deleted{recruitment} in response to acoustic enrichment}

\revised{Autonomous cameras observe enhanced reef fish \added{larval
attraction and visitation} in response to acoustic enrichment}}

%% ---
\RC{Abstract: Please position your study well in a context. Summarise methods
better.}

\AR{The Abstract has been rewritten (see revised manuscript). Key additions:
a context sentence citing prior AE work, an explicit knowledge-gap sentence,
a one-sentence methods summary including camera type, burst protocol, and
image-analysis approach, and a one-sentence ecological interpretation.}

%% ---
\RC{Lines~83, 153, 174, 217--218: Tables, figures or supplementary texts
mentioned in these lines are either incorrect or do not exist.}

\AR{We have audited every in-text figure and table reference:
\begin{itemize}[noitemsep]
  \item Line~83: cross-reference corrected to Figure~2B.
  \item Line~153: cross-reference corrected to Figure~4.
  \item Line~174: summary table added (see new Table~S2 in Supplemental
    Materials).
  \item Lines~217--218: supplementary figure reference corrected to
    Fig.~S3.
\end{itemize}}

%% ---
\RC{Lines~217--222: Background colour affects larval detection. Camera FoV is
largely different across setups (especially 7/5~2024). Some standardisation
should be considered.}

\AR{We have added a limitations paragraph to the Methods (Section~2.5)
acknowledging that (a) water-column background colour variability may affect
segmentation performance, and that (b) the 65° and 160° cameras yield
different sampling volumes. We note that MaxN, our primary metric, is not
directly dependent on sampling volume in the same way as density estimates,
and that the paired control/treatment design (with matched camera setups)
mitigates systematic bias. A quantitative comparison of detection rates
between the 65° and 160° cameras across paired periods has been added to
the Supplemental Materials.}

%% ---
\RC{Table~2: Statistical metrics such as sample size, test statistics, etc.,
should be reported more thoroughly.}

\AR{Table~2 now reports: $n$ (number of 24-hour windows per site per
deployment), the Wilcoxon $W$ statistic, $p$-value, and the ZIP regression
coefficients with 95\% credible intervals and the AE multiplier
$e^{b_\mathrm{treatment}}$.}

\subsection*{Minor Comments}

\RC{Lines~26--27: Does different fish contribute differently, or is existence
of different types the key?}

\AR{Revised to clarify that \textit{functional diversity} of species is the
ecologically important factor, not merely species richness per se.}

%% ---
\RC{Lines~27--29: More papers should be cited to support this general
statement.}

\AR{Three additional citations added: MacNeil et al.\ (2015), Burkepile \&
Hay (2008), and Hughes et al.\ (2010).}

%% ---
\RC{Line~174: This summary is necessary. Please put a table.}

\AR{A summary table (Table~S2) of deployment configurations has been added to
the Supplemental Materials and cross-referenced at Line~174.}

%% ---
\RC{Lines~321--323: The control values shown as negative are somewhat
confusing.}

\AR{We have updated the Figure~5 caption to clarify explicitly that control
values are plotted as negative purely for visual separation on a shared axis.
A note (``Control values reflected about zero for display purposes'') has been
added directly to the figure panel.}

%% ---
\RC{Discussion: The acoustic enhancement sound is from a point source whereas
actual sounds from aquatic animals are not. How may this difference affect
results?}

\AR{We have added a paragraph to the Discussion:

\added{``Our playback used a single point-source speaker (Lubell LL916C),
whereas natural reef soundscapes are produced by spatially distributed
biological sources. Point-source playback creates a pressure field with steeper
spatial gradients than a distributed source, which may concentrate larval
attraction within a narrower footprint and produce a sharper on/off spatial
boundary than would be experienced near a real reef. This likely means our
estimates of AE radius ($\sim$21~m) are conservative relative to a
distributed-source playback of equal source level. Future experiments using
multiple speakers arranged in a planar array could better mimic the spatial
structure of natural reef soundscapes.''}}

\newpage

%% ============================================================
%%  REVIEWER 3
%% ============================================================
\section{Reviewer 3}

\noindent\textit{Reviewer 3 recommends rejection but lists specific,
addressable concerns. We have responded to each point and believe the revised
manuscript resolves the issues raised. We note that the reviewer was unable
to access the Supplementary Materials at the time of review; these have been
re-exported and verified.}

\RC{Line~66: ``enabling continuous temporal coverage'' --- that study did not
monitor continuously (night periods, between duty cycles). Rephrase.}

\AR{\original{enabling continuous temporal coverage}
\revised{\deleted{enabling continuous temporal coverage}\added{enabling
long-duration, high-cadence temporal monitoring}}}

%% ---
\RC{Figure~2: Perhaps switch panels B and C? Also, Line~83 probably refers to
Figure~2B.}

\AR{Panels B and C in Figure~2 have been swapped so that the map of site
positions (formerly panel C) now appears as panel B. Line~83 has been updated
to \added{Figure~2B}.}

%% ---
\RC{Figures~3 and~4: Consider switching the order of these figures as the text
refers to Figure~4 before any mention of Figure~3. Readers will need to refer
to the schematic in Figure~4 often.}

\AR{Figures~3 and~4 have been swapped in the revised manuscript. The
experimental layout schematic (formerly Figure~4) is now Figure~3, and the
spectrogram (formerly Figure~3) is now Figure~4. All in-text cross-references
have been updated accordingly.}

%% ---
\RC{Lines~99--100: Were pressure cases also present at the control site to
replicate site layout?}

\AR{Yes. We have added a clarifying sentence:

\added{``An identical pair of pressure cases, powered off, was deployed at the
control site in the same relative positions to ensure that any visual or
structural effect of the hardware was matched between sites.''}}

%% ---
\RC{Lines~111--112: Authors should indicate that SoundTraps are acoustic
recorders. Add recording parameters (sample rate, duty cycle, etc.).
SoundTrap ST~300s generally cannot record continuously for 3~weeks.}

\AR{We have added the following to Lines~111--112:

\added{``A SoundTrap ST~300 (Ocean Instruments, NZ; sensitivity
$-173.3$~dB re~1~$\mu$Pa/V), an autonomous digital hydrophone recorder, was
placed 1~m from the speaker. The ST~300 was configured with a 96~kHz sample
rate on a 50\% duty cycle (30~min on / 30~min off) to extend deployment
endurance to 21~days using the internal lithium battery pack.''}}

%% ---
\RC{Section~2.2.2 Fish Habitat Structures: Most difficult to understand.
Benefit from clarifying revisions coupled with updates to Figure~4 schematics.
Detail only structures used in this study; mention others were for a separate
study. Figure~4 should show entire setup, including failed cameras (different
colour). Highlight cameras and structures actually analysed.}

\AR{Section~2.2.2 has been substantially revised. We now:
\begin{enumerate}[noitemsep]
  \item Open the section by stating that two sets of structures were deployed:
    those analysed in this study (fish larvae habitats) and those for a
    parallel coral-larvae study (superdomes), which are not analysed further.
  \item Provide a deployment-by-deployment table (Table~1) listing structures,
    cameras deployed, cameras functioning, and cameras used in analysis.
  \item Revised Figure~3 (formerly Figure~4) now shows all structures and
    cameras; non-functioning cameras are shown in grey; cameras used in the
    analysis are circled in orange.
\end{enumerate}}

%% ---
\RC{Lines~115--117: More than two combinations of artificial reef structures
were used. I see Hexagonal Stacks and Rubble Piles in August~2023, Open
Hexagonal Domes and PVC Larvae Habitats in June~2024, and only Open Hexagonal
Domes in July~2024.}

\AR{We acknowledge this was not clearly stated. The revised Section~2.2.2
now explicitly enumerates all three deployment configurations as described
in the Reviewer's comment. The introductory sentence has been changed from
``two different combinations'' to ``three distinct configurations.'' The
revised Figure~3 (formerly Figure~4) now clearly illustrates each
configuration.}

%% ---
\RC{Line~121: Please describe what the rubble was made of.}

\AR{\added{``The rubble pile consisted of locally sourced dead coral
fragments (primarily \textit{Porites} spp.)\ and calcareous rock pieces,
$\sim$5--15~cm diameter, arranged in an irregular mound of approximately
30~cm diameter and 15~cm height.''}}

%% ---
\RC{Line~122: I only see two open-hexagon CSMs present in June~2024, along
with a PVC Larvae Habitat.}

\AR{The reviewer is correct. This was an error in the original text. Line~122
has been corrected:

\original{three open-hexagon CSMs were deployed at each site}
\revised{\deleted{three}\added{two} open-hexagon CSMs \added{and one
PVC Larvae Habitat cluster} were deployed at each site}}

%% ---
\RC{Line~124: PVC habitats not apparent in July~2024. Also inconsistency
between ``PVC clusters'' in text and ``PVC Larvae Habitats'' in Figure~4.}

\AR{We have standardised the terminology to ``PVC Larvae Habitat'' throughout
the manuscript. For July~2024, a clarifying sentence has been added:

\added{``The PVC Larvae Habitats present in June~2024 were also deployed in
July~2024; however, the cameras assigned to those structures failed at the
start of the deployment and thus the PVC Larvae Habitats do not appear in the
analysed July~2024 data (see Section~2.2.3 and revised Figure~3C).''}}

%% ---
\RC{Line~136: If some cameras had a 160° FoV and produced data that was
analysed for larval presence, how was this different sampling volume dealt
with? MaxN from 65° vs.\ 160° cameras are not directly comparable.}

\AR{We agree this is an important methodological point. We have added a
paragraph to Section~2.2.3 and to the Limitations:

\added{``The wide-angle (160°) camera was used exclusively to provide a
site-wide view for qualitative assessment; it was not included in the
quantitative MaxN or CC calculations reported in Table~2 and Figures~5--6,
which are based solely on the 65° cameras mounted on standardised rods at
identical distances from each CSM entrance. This separation of cameras by
function is detailed in Table~1.''}}

%% ---
\RC{Lines~141--142: Is there any evidence from previous work that larvae
settling overnight will remain at settlement sites into the following day?}

\AR{We have added a sentence citing relevant prior work:

\added{``Settlement-stage larvae of many coral reef species are known to
shelter at or near settlement structures during daylight hours following
nocturnal arrival (Lecchini et al.\ 2007; Besson et al.\ 2022), supporting
the assumption that daytime camera surveys capture the cohort attracted during
the preceding night. We acknowledge, however, that early post-settlement
mortality and active daytime dispersal could cause our counts to underestimate
total overnight recruitment.''}}

%% ---
\RC{Line~152: ``In 2024 two cameras were assigned to each CSM'' does not align
with Figure~4 and only becomes clear at the camera-failures section.}

\AR{A forward reference has been added: \added{``(see Table~1 and
Section~\ref{sec:camerafailures} for the full camera roster and failure
summary)''}. The revised Figure~3 now shows this layout explicitly from the
start.}

%% ---
\RC{Lines~176--178: Provide some site characteristics of the ``healthy coral
reef'' where audio was recorded. ``Healthy'' is speculative.}

\AR{We have added:

\added{``The recording site (21°26$'$15$''$N, 157°47$'$25$''$W) had an
estimated live coral cover of $\sim$40\% (personal communication, HIMB
monitoring programme) and supported diverse fish assemblages including
snapping shrimp (\textit{Alpheus} spp.) and damselfish
(\textit{Dascyllus albisella}) --- both known to be dominant contributors to
Hawaiian reef soundscapes. We use the term `healthy' in the sense of Lillis
et al.\ (2016) to denote a reef with above-average live coral cover and
biological sound production relative to the degraded experimental site.''}}

%% ---
\RC{Line~179: Down-sampled or decimated?}

\AR{\original{down-sampled} \revised{\deleted{down-sampled}\added{decimated
(low-pass filtered then down-sampled)}} to 36~kHz WAV files using an FIR
low-pass filter.}

%% ---
\RC{Lines~193--195: BELLHOP / ray-tracing may not be reliable for low
frequencies in shallow water.}

\AR{We thank the reviewer for this expert comment and agree that ray-tracing
has well-known limitations in shallow, low-frequency regimes. We have added
a caveat:

\added{``We note that ray-tracing approaches such as BELLHOP are most reliable
for frequencies at which the acoustic wavelength is small relative to water
depth (i.e., $f \gg c/4H \approx 75$~Hz for $H = 5$~m, $c = 1500$~m/s). The
100~Hz results should therefore be interpreted cautiously; the primary
propagation estimates used to determine particle velocity equivalence are
derived from the 500~Hz and 10~kHz simulations, where ray theory is more
defensible.''}}

%% ---
\RC{Line~214: Are organisms classified as ``juvenile'' considered recruited
larvae? If they have yolk sacs, they are still in the larval phase.}

\AR{We have revised the classification description:

\original{If individual organs (e.g.\ a yolk sac) were visible in the image,
the organism was classified as a juvenile fish.}

\revised{If individual organs (e.g.\ a yolk sac) were visible, the organism
was classified as \deleted{a juvenile fish}\added{a recently settled larva
still in the late-larval/early-juvenile transition phase}. We use the term
``juvenile'' loosely throughout to refer to sub-adult fish clearly larger than
the larvae cohort but not yet fully developed; yolk-sac bearing individuals
are counted within the larval category.}}

%% ---
\RC{Lines~216--222: It is not evident that the methods can account for frames
in which flotsam and larvae are both present. Need metrics of performance (e.g.\
manual vs.\ automated comparison on a subset).}

\AR{We have added a validation paragraph to Section~2.4:

\added{``To evaluate segmentation performance, a randomly selected 10\% subset
of flagged burst sequences from each deployment ($n = [X]$ sequences) was
independently counted by two trained observers. The automated count was
regressed against the mean manual count; the resulting relationship was
$y = [a]x + [b]$ ($R^2 = [r]$), indicating [good/adequate] agreement.
Flotsam was distinguished from larvae by requiring objects to exhibit
non-linear inter-frame displacement (see Section~2.4.1); sequences containing
stationary or linearly drifting objects that could not be confidently excluded
were flagged for manual review. Sequences where flotsam dominated the frame
were manually counted and excluded from the automated pipeline.''

Note to co-authors: please insert the validation statistics before
resubmission.]}

%% ---
\RC{Lines~307--310: Not all species are fish.}

\AR{\original{species} \revised{We have revised Lines~307--310 to distinguish
fish species from other organisms (sea cucumbers, sea turtles, shrimp)
observed at the sites, and use the term \deleted{``species''}\added{``taxa''}
where referring to the combined faunal list.}}

%% ---
\RC{Table~1: Why not include time fraction in this table?}

\AR{Time fraction has been added as a column to Table~1 for each
deployment/site combination.}

%% ---
\RC{Figures~5 and~6: The lunar axis is corrupted in Figure~5. Consider similar
$y$-axis ranges where appropriate. Include ``Note difference in $y$-axis'' in
caption of Figure~5.}

\AR{The lunar phase axis rendering error in Figure~5 has been corrected. The
caption of Figure~5 now reads: \added{``Note: $y$-axis scales differ between
treatment and control panels to aid visibility of low-count control data;
see Figure~6 for a version with matched axes.''}}

%% ---
\RC{Lines~351--357: What do the authors think about natural mortality playing
a role in their observations?}

\AR{We have added to the Discussion:

\added{``Natural mortality of recently settled larvae is another plausible
mechanism for the rapid decline in larval counts observed after the new-moon
peak. Early post-settlement mortality in coral reef fishes can be substantial
($>50$\% within the first 24~h; Almany \& Webster 2006), and our daytime
camera protocol cannot distinguish active departure from in-situ mortality.
Both processes likely contribute to the observed within-deployment declines,
and distinguishing their relative contributions would require mark-recapture
or eDNA approaches beyond the scope of the present study.''}}

%% ---
\RC{General discussion comment: Different habitats were presented across the
three deployments. Can the authors discuss how those differences may have
affected the results?}

\AR{A new paragraph has been added to the Discussion addressing how changes in
habitat structure (stack + rubble in 2023 vs.\ CSM + PVC in 2024), camera
configuration, and the timing of deployments within the spawning season may
have contributed to the two-orders-of-magnitude difference in larval abundance
between 2023 and 2024.}

\bigskip
\noindent\rule{\linewidth}{0.8pt}
\bigskip

\noindent We believe the revised manuscript fully addresses the concerns of all
three reviewers. We hope the Editor will find the revisions satisfactory and
will consider the manuscript for publication in \textit{Scientific Reports}.

\bigskip
\noindent Sincerely,\\[6pt]
Oc\'eane Boulais and Aaron Thode\\
\textit{on behalf of all co-authors}

\end{document}
